\section*{Chapter \ref{ch:m2:emerging-market-fx-and-ndfs} --- Emerging-Market FX and NDFs}

\begin{solution}{exo:m2:emerging-market-fx-and-ndfs:1}
The opposite of the buyer: it pays USD~319\,186 at 1,425.50 and receives
USD~222\,222 at 1,350.00.
\end{solution}

\begin{solution}{exo:m2:emerging-market-fx-and-ndfs:2}
Rupees cannot be delivered offshore, and onshore hedging is restricted by the
controls on who may trade and why. The NDF settles in dollars outside India
against the published fixing, so the fund can hedge without moving rupees.
\end{solution}

\begin{solution}{exo:m2:emerging-market-fx-and-ndfs:3}
6.9580 to 7.2420, 2\% either side of 7.1000; the Hong Kong dollar, 7.75 to 7.85
per US dollar.
\end{solution}

\begin{solution}{exo:m2:emerging-market-fx-and-ndfs:4}
Two business days before Monday 5 October, skipping the weekend and the Friday
holiday: Wednesday 30 September 2026.
\end{solution}

\begin{solution}{exo:m2:emerging-market-fx-and-ndfs:5}
Offshore forwards imply a higher renminbi rate than onshore ones. An investor
that hedges renminbi bonds by selling renminbi forward earns the dollar rate less
the implied renminbi rate on the hedge: relative to spot, the renminbi's forward
premium is smaller offshore (7.12 to 7.07) than onshore (7.10 to 7.04), so the
offshore hedge returns fewer dollars. It costs about the basis more, and \omterm{def:m2:emerging-market-fx-and-ndfs:controls}{capital
controls} keep the investor from using the onshore price.
\end{solution}

\begin{solution}{exo:m2:emerging-market-fx-and-ndfs:6}
To stop its currency rising the central bank sells its own currency, which it
can create without limit; it accumulates foreign reserves, a risk to its balance
sheet but not a constraint. To stop its currency falling it must sell foreign
currency, which runs out.
\end{solution}

\begin{solution}{exo:m2:emerging-market-fx-and-ndfs:7}
The lowest reference rate while the floor held was 1.2008, on 1 June 2012. It was
1.2010 on 14 January 2015 and 1.0280 on 15 January; the lowest in January was
0.9816, on 23 January.
\end{solution}

\begin{solution}{exo:m2:emerging-market-fx-and-ndfs:8}
The low volatility measures the market's belief in the floor, not the size of
the move if it goes; the risk is a jump, which a stop does not limit. A stop
becomes a market order and fills at the next price, which on 15 January 2015 was
more than 14\% lower by the afternoon reference. The position's risk was the
floor's removal times the jump, not 2\%.
\end{solution}

\begin{solution}{pb:m2:emerging-market-fx-and-ndfs:1}
\textbf{1.} $1\,000\,000 \times 1.2010/20 = $ CHF~60\,050.
\textbf{2.} $1\,000\,000 \times (1.2010 - 1.19) = $ CHF~11\,000, had it filled at
its level.
\textbf{3.} $1 - 1.0280/1.2010 = 14.4\%$.
\textbf{4.} CHF~173\,000, 2.88 times the margin.
\textbf{5.} $-$CHF~112\,950: the client owes the broker that amount.
\textbf{6.} An order to sell once the price trades at or below its level; it then
becomes a market order and guarantees execution, not the price.
\textbf{7.} Everyone who had relied on the floor wanted to sell euros at once, and
buyers stepped away until they could judge where the rate would settle.
\textbf{8.} Years of trading just above 1.20 made the downside look capped; the
leverage was sized for the observed volatility, not for the policy's end.
\textbf{9.} Low, reflecting the floor's credibility; it missed the size of the
jump if the floor went, because that risk is not continuous.
\textbf{10.} A put struck near 1.20 was cheap before and paid about 17 centimes
per euro that afternoon: it would have covered the loss where the stop failed.
\textbf{11.} The broker, which owes its own counterparties for the client's
losing positions whether or not the client pays.
\textbf{12.} With a USD~300 million two-year loan from Leucadia, agreed on 16
January 2015, whose proceeds replaced the capital of its regulated entities.
\textbf{13.} An initial 10\% a year, rising by 1.5 points a quarter to at most
17\%, with covenants and mandatory prepayments.
\textbf{14.} By the jump the policy's end could cause, not by recent volatility,
and with margin that covers it.
\textbf{15.} Negative-balance protection, where a regulator requires it of
retail brokers for leveraged products: client losses are capped at the deposit.
\textbf{16.} A gradual exit would be anticipated and attacked; a surprise avoids
buying unlimited foreign currency in the days before, at the cost of a jump.
\textbf{17.} Any one-sided commitment: other currency floors and ceilings, pegs
defended with limited reserves, and bands managed close to their limits.
\textbf{18.} By the loss if the peg breaks, sized from the jump and the capital
one can lose, not from the pegged rate's volatility.
\textbf{19.} Named result: \emph{the floor that broke} cost the 20-to-1 long
2.88 times its margin, CHF~173\,000 on EUR~1 million at the ECB reference rate of
15 January 2015, leaving the client CHF~112\,950 in debt.
\textbf{20.} It is calm until the policy ends, and then it jumps.
\end{solution}

\begin{solution}{iq:m2:emerging-market-fx-and-ndfs:1}
A forward on a currency pair settled only in dollars, for the difference between
the contract rate and a fixing published two business days before settlement:
the dollar buyer receives $N(S-K)/S$. It exists for currencies that cannot be
delivered offshore.
\iqlookfor{cash settlement against a named fixing, and why it exists.}
\end{solution}

\begin{solution}{iq:m2:emerging-market-fx-and-ndfs:2}
\omterm{def:m2:emerging-market-fx-and-ndfs:controls}{Capital controls} separate the markets: money cannot move freely between
mainland China and Hong Kong, so onshore and offshore supply and demand can
differ, and arbitrage through the permitted channels is limited in size. The
difference widens when flows are one-sided and the channels bind.
\iqlookfor{the controls and the limits to arbitrage.}
\end{solution}

\begin{solution}{iq:m2:emerging-market-fx-and-ndfs:3}
As a mixture: a diffusion around the peg with low volatility, plus a jump of
uncertain size at a random time whose intensity reflects the policy's
credibility (a jump-diffusion). Calibrate the jump intensity and size to risk
reversals and far out-of-the-money options, which price the break; the at-the-money
volatility says little.
\iqlookfor{a jump component, and where its parameters come from.}
\end{solution}

\begin{solution}{iq:m2:emerging-market-fx-and-ndfs:4}
Positions sized for a calm, capped market met a 14\% jump; stops filled far below
their levels; leveraged clients lost more than their deposits, and brokers who
could not collect the debit balances had to raise capital or failed. Risk models
had used recent volatility rather than the event.
\iqlookfor{jump risk, failed stops, leverage and counterparty loss.}
\end{solution}

\begin{solution}{iq:m2:emerging-market-fx-and-ndfs:5}
Sell reais forward against dollars with NDFs against the PTAX fixing, rolled or
matched to the bonds' horizon; check onshore versus offshore pricing and the
basis. At the fixing: publication delays or disruptions (the standard fallbacks
apply), a fixing that differs from the rate at which the bonds are sold, and
the basis moving against the hedge.
\iqlookfor{NDF against the named fixing, basis, and fixing risk.}
\end{solution}

\begin{solution}{iq:m2:emerging-market-fx-and-ndfs:6}
Margin by stress scenarios as well as by volatility: jumps on pegs and after
known events, applied pair by pair; leverage caps per client and pair;
concentration limits on crowded positions; hedge the aggregate client book
with liquidity providers and model the provider's own gaps; capital sized to
the worst scenario's uncollectable debit balances; negative-balance protection
priced in; real-time monitoring with automatic de-risking before scheduled
events.
\iqlookfor{scenario margin, leverage caps and capital for gap risk.}
\end{solution}
